\chapter{The Number Splits Into Faces}

\section{One corridor, four faces}

The machine owns a reading: a bracket, floored, reproducible. This part asks what such a reading \emph{is}, and the
first answer is that it is not one thing but four. The same bracket, read four ways, presents four faces --- and the four
faces are the four gauges of the whole series. This section builds the claim carefully, because it is the structural
reason there are four books and not one: a single reading looks like four different quantities depending on which
apparatus is turned on it.

Picture the reading as a \emph{corridor} rather than a point --- a bracketed residue carried down the machine's descent,
a channel along which one quantity travels. A channel can be read at more than one port. Shannon's~\cite{shannon1948}
account of a signal already contains the idea: one transmission supports many read-outs, and what you recover depends on
the detector you attach. The machine's corridor is like this. Attach one apparatus and the bracket reads as a
\emph{count} --- the tally of turns the seam accumulated, a whole number of passes. Attach another and it reads as a
\emph{bound} --- the incompressible width the descent could not close, a magnitude the residue resists shrinking below.
Attach a third and it reads as a \emph{decided bit} --- an orientation, which of two ways a carried fact falls. Attach a
fourth and it reads as a \emph{cost} --- the price in the machine's own currency of taking the reading at all. One
corridor, four ports, four faces.

These four faces are the four gauges' native quantities, and it is worth naming which is which so the partition is
clear. A physical gauge, reading the same corridor, calls the count a \emph{charge}, the bound a \emph{mass}, the
decided bit a \emph{phase}, the cost a \emph{value} or an action; those are that gauge's names, and this book neither
builds nor defends them. This gauge keeps the computational names: the count, the bound, the bit, the cost. The point is
not that one set of names is right and the other wrong. The point is that they are the \emph{same corridor}, read at
different ports, and the faces differ only in which apparatus is attached. The residue does not change from face to
face; only the reading does.

In the two verbs the four faces are four funges of one corridor. Each apparatus funges the bracket into a single
quantity of its own kind --- the counting apparatus pools it into a tally, the bounding apparatus into an interval, the
deciding apparatus into a bit, the pricing apparatus into a cost. And each face is tanged apart from the others by which
port read it, not by any difference in the corridor itself. So the number splits into faces the way one signal splits
into readings: one thing tanged four ways by four detectors, funged four times into four quantities, all of them
readings of the single residue running down the corridor.

\section{The partition point}

If one corridor carries four faces, then somewhere there is a point at which the four gauges divide the labor --- a place
where the single machine is partitioned into four readings, none of which owns the whole. This section names that
point, because it is the structural seam of the entire series. No gauge reads the whole machine alone; each reads one
face exactly and hands the others off. The partition point is where that division is made, and understanding it is
understanding why four decorrelated books are stronger than one.

The division is not a carving of the machine into four parts. It is a choice of \emph{which port to read} --- and the
four gauges make four different choices, each reading the corridor at the port where its own language is exact. This
gauge reads at the ports where the corridor speaks computation: the count, the cost, the bound as an incompressibility,
the bit as a decision. It is precise there and vague elsewhere; asked to read the corridor as a physical quantity, it
has nothing to say, because it carries no such apparatus. A physical gauge is the mirror image: exact at the ports where
the corridor reads as force or field, mute at the ports where it reads as elaboration cost. The partition point is where
each gauge commits to its own ports and cedes the rest.

The reason the partition strengthens rather than weakens the reading is the deepest structural fact in the series, and
it is worth stating exactly. A single gauge, reading the corridor at all four ports with one apparatus, could tune its
readings to agree with each other --- could adjust its own machinery until the four faces lined up however it wished.
Four gauges that share no apparatus cannot do this. Each reads its own port in its own language, blind to what the
others read, and none can adjust its reading to flatter the rest. So when the four faces turn out to be readings of
\emph{one} corridor --- when the count, the bound, the bit, and the cost are found to be one residue seen four ways ---
that agreement is not a tuning but a fact. The partition decorrelates the gauges; decorrelation is what makes their
eventual agreement mean something. A truth four independent instruments cannot see each other reach is a truth none of
them manufactured.

There is a caution the partition builds in, and this book keeps it strictly. Because no gauge owns the whole machine,
no gauge may claim the whole machine's reading. This volume reads the computational faces and says so; it does not reach
across the partition to assert the physical ones, and where it names them it names them as another gauge's, marked and
handed off. The partition is a discipline of modesty as much as a division of labor: each gauge speaks only for the
ports it can read, and the four together speak for the corridor precisely because none of them overreaches into the
others' ports. The strength of the four-book instrument is exactly this refusal of any one book to read the whole.

In the two verbs the partition point is where one corridor is tanged into four readings by four apparatus and then, at
the very end, funged back into one. Each gauge tanges its own face apart --- selects the port its language reads --- and
works that face alone. The series funges the four faces back together only at the close, when the four gauges' brackets
are set side by side and found to pool into one interval. The partition tanges the corridor apart into faces so that the
ending can funge them back into a single trusted bound. That the four faces are one residue is the next section's, and
it is the fact the whole partition exists to earn.

\section{Same residue, different reading}

The partition splits the corridor into four faces, read by four gauges. This closing section holds the four faces up
against each other and states the fact that makes them worth splitting: they are the \emph{same residue}. Not four
residues, one per gauge --- one residue, read four ways. The face depends entirely on where the corridor is viewed, and
on nothing about the residue itself. This is the claim the whole series rests on, and this gauge states it plainly for
its own two faces before the rest of the book reads them out.

Take the count and the bound, the two faces this gauge reads most directly, and see that they are one thing. The count
is the tally of turns the descent accumulated; the bound is the width the descent could not close. But the turns and the
width are not independent quantities that happen to travel together --- they are the same descent, read once by its
length and once by its remainder. A descent that ran more turns closed more of the residue and left a narrower bound; a
descent that ran fewer left a wider one. The count and the bound co-vary because they are two readings of a single walk
down the tower, one reading its length and the other its leftover. They look like different quantities only because they
are read at different ports; underneath, they are the one residue the descent carried.

The face depends on the viewing angle, and the machine can even say what changes the angle. Read the corridor by
\emph{how far the descent went} and you get the count. Read it by \emph{what the descent could not resolve} and you get
the bound. Read it by \emph{which way a carried fact fell} and you get the bit. Read it by \emph{what the descent cost}
and you get the cost. The residue never stirred; the reader did --- attaching first one apparatus, then another, to the
same corridor. This is why the book can be strict about the residue and generous about its readings: there is one thing
to be honest about, the residue, and many honest ways to read it. The discipline is on the residue; the readings are
where the four gauges are free to differ.

And this is the fact that arms the ending. If the four faces were four different quantities, four gauges agreeing on a
bracket would be four coincidences. But the four faces are one residue, so four gauges agreeing on a bracket is four
honest readings of a single thing --- and a single thing read four incompatible ways, all landing on one interval, is a
thing whose interval is real. The book plants the claim here, at the partition, so that when the far chapters read the
count, the bound, the bit, and the cost each in turn, the reader already knows they are reading one residue, and that
their agreement at the end is not luck but the residue showing the same width through four different windows.

In the two verbs the identity of the four faces is the funge the series is built to close. Four gauges tange the corridor
into four faces and read each apart. But the residue underneath is one, and so the four readings funge back together ---
pool, at the end, into the single thing they were always readings of. The partition's tanging is in service of this final
funge: tell the corridor apart into faces so thoroughly, in four languages that cannot see each other, that when the
faces pool back into one bracket the pooling cannot be doubted. Same residue, different reading --- and the next chapter
begins reading the faces out, starting where the corridor first settles into a stable, nameable object: the electron.
